% from ega4-2-fr.tex
\oldpage[IV-2]{71}

donc, pour tout point $x'$ de $X'=X\otimes_k k'$ au-dessus de $x$,
$\mathcal O_{X',x'}$, qui est isomorphe à un anneau local de
$\operatorname{Spec}(\mathcal O_{X,x})\otimes_k k'$, est réduit (resp.
intègre).
\end{proof}

\begin{proposition}[4.6.12]
\label{IV.2.4.6.12-fr}
Soient $k$ un corps~; $X$ un $k$-préschéma, $x$ un point de $X$, $\Omega$ une
extension parfaite de $k$. Les conditions suivantes sont équivalentes~:
\begin{enumerate}
\item[a)] $X$ est géométriquement réduit sur $k$ au point $x$, autrement dit,
pour toute extension $k'$ de $k$ et tout point $x'$ de
$X'=X\otimes_k k'$ au-dessus de $x$, $\mathcal O_{X',x'}$ est réduit.
\item[b)] Le préschéma $X\otimes_k\Omega$ est réduit en un point au-dessus de
$x$.
\item[c)] Pour toute extension finie radicielle $k'$ de $k$,
$X'=X\otimes_k k'$ est réduit en l'unique point au-dessus de $x$.
\item[d)] $\operatorname{Spec}(\mathcal O_{X,x})$ est géométriquement réduit
sur $k$.
\item[e)] $\mathcal O_{X,x}$ est réduit, et pour toute composante irréductible
$Z$ de $X$ contenant $x$, de point générique $z$, $k(z)$ est une extension
séparable de $k$.
\end{enumerate}
\end{proposition}

\begin{proof}
Les implications d) $\Rightarrow$ a) $\Rightarrow$ b) $\Rightarrow$ c) sont
immédiates, en tenant compte pour la dernière de (4.6.10), (4.6.11) et de ce
que toute extension radicielle de $k$ est isomorphe à une sous-extension de
$\Omega$. Prouvons que c) implique d). Il suffit de montrer que pour toute
extension radicielle finie $k'$ de $k$, $\mathcal O_{X,x}\otimes_k k'$ est
réduit (4.6.1)~; or cet anneau est l'anneau local de $X'=X\otimes_k k'$ en
l'unique point $x'$ de $X'$ au-dessus de $x$ (I, 3.5.8 et 3.5.7)~; il est donc
réduit en vertu de l'hypothèse c). Enfin, l'équivalence de d) et de e)
résulte de l'équivalence de b) et e) dans (4.6.1), puisque les composantes
irréductibles de $X$ contenant $x$ correspondent biunivoquement aux composantes
irréductibles de $\operatorname{Spec}(\mathcal O_{X,x})$ (I, 2.4.2).
\end{proof}

\begin{corollary}[4.6.13]
\label{IV.2.4.6.13-fr}
Sous les hypothèses de (4.6.12), supposons de plus que $X$ soit localement
noethérien. Alors les conditions a) à e) de (4.6.12) équivalent aussi à la
suivante~:
\begin{enumerate}
\item[f)] Il existe un voisinage ouvert $U$ de $x$ qui est géométriquement
réduit sur $k$.
\end{enumerate}
\end{corollary}

\begin{proof}
En effet, cette condition implique trivialement que $X$ est géométriquement
réduit au point $x$ sur $k$. Inversement, s'il en est ainsi, puisque l'anneau
$\mathcal O_{X,x}$ est réduit, il existe un voisinage ouvert $U$ de $x$ dans
$X$ qui est réduit (I, 6.1.13)~; prenant $U$ assez petit, on peut en outre
supposer que $U$ ne rencontre aucune des composantes irréductibles de $X$ ne
contenant pas $x$. Alors le critère (4.6.12, e)) prouve que $U$ est
géométriquement réduit sur $k$, compte tenu du critère (4.6.1, e)).
\end{proof}

Il résulte de (4.6.13) que lorsque $X$ est localement noethérien, l'ensemble
des $x\in X$ où $X$ est géométriquement réduit sur $k$ est \emph{ouvert}, et
c'est le plus grand ouvert de $X$ qui soit géométriquement réduit sur $k$.

\begin{env}[4.6.14]
\label{IV.2.4.6.14-fr}
Il résulte de (4.6.10) et (4.6.11) que si $\Omega$ est une extension
algébriquement close de $k$, il revient au même de dire que $X$ est
géométriquement ponctuellement intègre au point $x$, ou que $X\otimes_k\Omega$
est intègre en un point au-dessus de $x$.

Pour que $X$ soit géométriquement ponctuellement intègre au point $x$, il est
nécessaire que $X$ soit intègre au point $x$, et que si $z$ est le point
générique de l'unique composante irréductible de $X$ contenant $x$, $k(z)$ soit
une extension séparable de $k$~; cela résulte de (4.6.10). Mais ces conditions
\emph{ne sont pas suffisantes}, comme il résulte de l'exemple (4.5.12, (ii)).
Une condition suffisante mais non nécessaire est que $X$ soit géométriquement
réduit au point $x$ et n'appartienne qu'à une seule composante irréductible de
$X$, que

